\documentclass[10pt,a4paper]{amsart}
\usepackage[margin=19mm]{geometry}
\usepackage{amsmath,amssymb,mathtools}
\usepackage[scr=rsfs,cal=euler]{mathalfa}
\usepackage{xfrac}
\usepackage[unicode,hidelinks,bookmarks=false]{hyperref}
\newcommand{\ie}{i.e.\ }
\newcommand{\cf}{cf.\ }
\newcommand{\resp}{respectively\ }
\newcommand{\Subsection}{\subsection}
\providecommand{\nobreakdash}{\nobreak\hbox{-}}
\setlength{\emergencystretch}{2em}
\multlinegap0pt
\usepackage{amssymb}
\newtheorem{proposition}{Proposition}
\newcommand{\J}{\mathrm{J}}
\providecommand{\llave}[1]{\left\lbrace #1\right\rbrace }
\providecommand{\norm}[1]{\left\lVert#1\right\rVert}
\providecommand{\pare}[1]{\left( #1\right) }
\newcommand{\tien}{\rightarrow}
\title{$B$-Multiplier Spaces}
\author{Rafael Correa-Morales, Fernando Galaz-Fontes}
\date{Source: arXiv:2608.13837v1 (CC BY 4.0)}
\begin{document}
\maketitle

\noindent\emph{Source and licence.} $B$-Multiplier Spaces. Version arXiv:2608.13837v1 (CC BY 4.0), distributed by arXiv under the Creative Commons Attribution 4.0 International licence, http://creativecommons.org/licenses/by/4.0/, as recorded in the arXiv licence metadata for this exact version and shown on the abstract page https://arxiv.org/abs/2608.13837v1. Full licence notice: \texttt{licenses/arxiv-2608.13837-CC-BY-4.0.txt}.\par\smallskip

\noindent\textit{Editorial scope.} Counted unit: the article's proposition final7 (source0.tex lines 1491--1494). The article's complete original proof of it is delivered with it (source0.tex lines 1495--1558, one \texttt{proof} environment of 64 lines). No mathematics is written, repaired or re-derived: the statement and the proof are the article's own lines, in the article's own notation, and no retained line is substituted or re-flowed.  No further source-local context is needed: every cross-reference of every retained line resolves inside the two delivered spans.  Nothing was omitted from any retained span, so the package carries no omission record.  No retained line contains a citation, so the wrapper carries no bibliography and no bibliography entry.  No retained line contains a figure environment, an image-inclusion command or drawing code, so no image is delivered and the package declares no asset dependency.  Editorial formatting only, in addition: an AMS \texttt{amsart} preamble that restates the article's own declarations and macros used by the retained lines, namely the environments \texttt{proposition} on separate counters, together with the standard packages those lines need.  The retained environments print in this document as Proposition 1 in order of appearance, whereas the article prints the same items with the numbering of its own full document.  The complete native source of the article as distributed in the arXiv e-print for this exact version is bundled unchanged as \texttt{sources/207617/source0.tex}.
 \begin{proposition}\label{final7}
 	Let $U \subseteq F(D,X)$ be a Banach space that is also a v-ideal. Then $U$ has an equivalent norm under which $U$ is a Banach v-ideal if, and only if,
 	$k:=\sup \llave{ \frac{ \|f\|}{\|g\|}:   f, g  \in U, J_f \leq Jg } < \infty$.
 \end{proposition}

 \begin{proof}
First, suppose that $U$ has a an equivalent norm   $\norm{\cdot}_{v}$ with respect to which $U$ is a Banach v-ideal. Then there exist $c,d\in(0,\infty)$ such that
$
c\norm{u}\leq \norm{u}_{v}\leq d\norm{u},  \forall u\in U.
$
Let $f,g\in U$ be such that $\J f\leq \J g$. Then
$
c\norm{f}\leq \norm{f}_{v}\leq \norm{g}_{v}\leq d\norm{g},
$
and therefore
$
k <\infty.
$ 

Now suppose that $k<\infty$. Consider the function $\norm{\cdot}_{v}:U\tien[0,\infty]$ defined by
\begin{equation}
	\norm{u}_{v}:=\sup\llave{\norm{h}:Jh\leq Ju}.
\end{equation}
To see that $\norm{\cdot}_{v}$ is a norm on $U$, we only prove the triangle inequality. Let $f,g,h\in U$ be such that $Jh\leq J\pare{f+g}$.
Given $d\in D$, it follows that
$
\norm{h\pare{d}}
\leq
\norm{f\pare{d}+g\pare{d}}
\leq
\norm{f\pare{d}}+\norm{g\pare{d}}.
$
Now define
\begin{equation}
	h_{1}\pare{d}:=
	\frac{\norm{f\pare{d}}h\pare{d}}{\norm{f\pare{d}}+\norm{g\pare{d}}},
	\qquad
	h_{2}\pare{d}:=
	\frac{\norm{g\pare{d}}h\pare{d}}{\norm{f\pare{d}}+\norm{g\pare{d}}},
	\qquad
	\text{if }\norm{f\pare{d}}+\norm{g\pare{d}}>0,
\end{equation}
and let $h_{1}(d)=h_{2}(d)=0$ whenever
$
\norm{f\pare{d}}+\norm{g\pare{d}}=0.
$
Consider the functions $h_{1}:D\tien X$ and $h_{2}:D\tien X$. Clearly, $Jh_{1}\leq Jf$, $Jh_{2}\leq Jg$, and $h=h_{1}+h_{2}$. Since $U$ is a $v$-ideal, it follows that $h_{1}, h_{2}\in U$. Therefore,
\begin{equation}
	\norm{h}=\norm{h_{1}+h_{2}}\leq \norm{h_{1}}+\norm{h_{2}}\leq \norm{f}_{v}+\norm{g}_{v},
\end{equation}
and it follows that 
$
\norm{f+g}_{v}\leq \norm{f}_{v}+\norm{g}_{v}.
$
Clearly
\begin{equation}\label{yaya6}
	\norm{u}\leq \norm{u}_{v},\qquad \forall u\in U.
\end{equation}
Moreover, if $f,g\in U$ satisfy $Jf\leq Jg$, then $\norm{f}_{v}\leq \norm{g}_{v}$. To prove that $\norm{\cdot}_{v}$ is complete, it is enough to show that it is equivalent to the original norm on $U$. 
Let $f,g\in U$ be such that $Jg\leq Jf$. Therefore
$
\norm{g}\leq k\norm{f}.
$
Taking the supremum over such  $g$, we obtain
$
\norm{f}_{v}\leq k\norm{f}.
$
Together with (\ref{yaya6}), this shows that the norm on $U$ and $\norm{\cdot}_{v}$ are equivalent.
 \end{proof}

\end{document}
